\section{Propagation from a characteristic-two curve}
\label{sec:propagation}

The input to this section is a pro-modular curve whose generic
representation is not virtually solvable.  We prove that completed forms
have full support on the deformation germ at that curve.  The residue
field of the germ is a Laurent-series field; accordingly, the patching
argument below uses finite jets, rather than compactness of a finite
residue field.

The argument follows the auxiliary-level framework of
Taylor--Wiles~\cite{TaylorWiles1995}, its formulation relative to local
rings by Kisin~\cite[Section~3.3]{Kisin2009FiniteFlat}, and the
characteristic-two determinant-twist construction of
Khare--Wintenberger~\cite[Section~9.1.3]{KhareWintenberger2009II}.
Pan's localized propagation argument is the global predecessor
\cite{Pan2022}. Here the curve residue field requires explicit control
of finite jets and fixed old-level quotients. We first construct
finite auxiliary stages, then preserve their filtration in the limit,
and finally use twists, sign intertwiners and diamond augmentation to
recover support with the original deformation action.

\begin{theorem}[Localized propagation]\label{thm:localized-propagation}
Let \(F\) be a totally real solvable field of even degree, completely
split at~\(2\), and let \(S\) contain its real and dyadic places and
the allowed finite ramification. Fix the determinant
\(\chi=\det r=\chi_{\mathrm{fin}}\eps^w\) of a target \(r\) as in
Section~\ref{sec:preliminaries}, a semisimple residual representation
with determinant \(\bar\chi\), and let
\(R_F^{\mathrm{ps},\chi}\) be the corresponding complete local global
pseudodeformation ring. Let $C$ be a characteristic-two curve in
\(\Spec R_F^{\mathrm{ps},\chi}\).  Suppose that $C$ is potentially pro-modular,
that its generic representation is not virtually solvable, and that its
local images at the non-dyadic finite places in $S$ are finite.  Every
irreducible closed locus of \(\Spec R_F^{\mathrm{ps},\chi}\)
containing $C$ is potentially pro-modular, with ramification allowed
only at the specified set \(S\) before base change.
\end{theorem}

We will prove the theorem at the end of the section.  First make a
totally real solvable Galois base change, split completely above $2$, on which
the curve is pro-modular.  Lemma~\ref{lem:inertia-basechange}, applied
both to the source deformation rings and to the Hecke algebras, permits
a further such base change with the following properties: the curve
representation is trivial at every bad finite place; the determinant
$\chi$ is unramified there; the source has unipotent inertia there; and
each bad residue cardinality is $1$ modulo $4$. By
Lemma~\ref{lem:forward-basechange}, the curve is then in
the support of forms with Iwahori level at the bad places and
hyperspecial level elsewhere away from $2$.  The assertion about the
source is set-theoretic, which is sufficient for the theorem.

Retain the notation $F,S$ after these base changes, and put
\[
 d=[F:\QQ],\qquad
 u=\#\{v\in S:v\nmid 2\infty\},\qquad s=2d+u.
\]
Use Lemma~\ref{lem:curve-completion} to enhance the curve by $z$ and
complete.  Its coefficient ring and residue field are
\[
 \mathcal L=\widehat{\OO[[z]]_{(\varpi)}},\qquad K=\kk((t)),
 \qquad z\longmapsto t.
\]
Finite extensions of these coefficients are allowed throughout.  Write
$\rho_*:G_{F,S}\to\GL_2(K)$ for the curve representation.  It is
absolutely irreducible, and remains so under the base changes just made.
All local and global deformation rings in this section are completed
at this same enhanced curve and have residue field $K$.

\subsection{Auxiliary primes and characters}

Character twists will enlarge the support obtained from completed
forms. Since twisting multiplies the determinant by the square of the
character, we need a variable-determinant problem as well as the
fixed-determinant one. In the former, the determinant is fixed at
the dyadic and real
places and on bad inertia, but is allowed to vary globally.  Enlarge
$\OO$ so that $i^2=-1$ belongs to it.  At a bad place, choose a tame
inertia generator $\tau$ and an arithmetic Frobenius $\phi$ with
$\phi\tau\phi^{-1}=\tau^{q_v}$.  For $\eta\in\{1,i\}$ impose
\[
 \det\rho(\tau)=1,\qquad
 \tr\rho(\tau)=\eta+\eta^{-1}.
\]
The determinant of $\rho(\phi)$ is unrestricted in the variable
problem.  The local curve is trivial, so wild inertia and prime-to-two
tame inertia disappear in these completed problems.  Denote the
completed product of the framed local rings at $S$ by
$R_{\mathrm{loc},\eta}$.
The $\eta=1$ condition contains the source with unipotent bad inertia.
The $\eta=i$ condition provides an irreducible comparison model:
Proposition~\ref{prop:local-geometry} identifies its special fiber
with that of the $1$-model and specifies the horizontal components
through each generic special-fiber point. We will use this comparison
to transfer support to every component of the $1$-model.

For a set $Q_n$ of auxiliary primes, let $R_{\eta,n}$ be the
$S$-framed global ring with these local conditions and unrestricted
additional ramification at $Q_n$.  A superscript $f$ denotes the
quotient imposing the global determinant $\chi$.  Write $R_\eta$ and
$R_\eta^f$ for the corresponding rings without $Q_n$.  Thus forgetting
the additional ramification gives quotient maps
\[
 R_{\eta,n}\twoheadrightarrow R_\eta,
 \qquad R_{\eta,n}^f\twoheadrightarrow R_\eta^f.
\]
The deformation interpretation over $\mathcal L$ and noetherianity
follow from Lemma~\ref{lem:curve-completion} and the relative tangent
and obstruction calculation of Proposition~\ref{prop:relative-presentation}.

\begin{lemma}[Auxiliary-prime bound]\label{lem:auxiliary-primes}
There is a sufficiently large integer $h$, independent of $n$, and
sets $Q_n$ of $h$ primes outside $S$ such that:
\begin{enumerate}
\item every prime in $Q_n$ splits in
$F(\zeta_{2^n},2^{1/2^n})$ and has residue cardinality $1$ modulo $2^n$;
\item its $\rho_*$-Frobenius is regular semisimple, with both eigenvalues
in one fixed finite extension of $K$;
\item there are surjections
\begin{equation}\label{eq:patching-presentation}
 B_\eta:=R_{\mathrm{loc},\eta}[[X_1,\ldots,X_g]]
 \twoheadrightarrow R_{\eta,n},\qquad g=2h+1.
\end{equation}
\end{enumerate}
The two sets of maps in \eqref{eq:patching-presentation}, for $\eta=1,i$,
can be chosen to agree modulo $\varpi$. For each \(\eta\), their
composites with \(R_{\eta,n}\twoheadrightarrow R_\eta\) can also be
chosen independently of \(n\).
\end{lemma}

\begin{proof}
Put $M=\ad\rho_*$.  The pairing $(A,B)\mapsto\tr(AB)$ on $M_2(K)$
is perfect also in characteristic two.  Thus the dual coefficient
module is again $M$: the mod-two cyclotomic twist is trivial.  The
relative Selmer condition is zero at $S$ and full at $Q_n$.  The dual
condition therefore consists of classes unramified outside $S$ whose
localizations at $Q_n$ vanish.

Set
\[
 F_\infty=F(\zeta_{2^\infty},2^{1/2^\infty}),\qquad
 H=\ker(G_{F,S}\longrightarrow\Gal(F_\infty/F)).
\]
The tower group is solvable.  Nonvirtual solvability of $\rho_*$
therefore implies that the geometric reduced Zariski closure of
$\rho_*(H)$ contains $\SL_2$; the same is true for every open subgroup
of $H$. Indeed, take the closure over an algebraic closure of \(K\).
Its smooth connected derived group lies in \(\SL_2\). If proper,
it has dimension at most two and is solvable, forcing the identity
component itself to be solvable, a contradiction
\cite[Corollaries~6.19 and~17.27]{MilneAlgebraicGroups}.
Consequently regular semisimple elements are dense in every open
coset of $H$.

We give the detection argument explicitly.
Let $c\in Z^1(G_{F,S},M)$ have zero diagonal evaluation at every
regular semisimple $\sigma\in H$.  Then the first-order
representation
\[
 \rho_c(\sigma)=(1+\epsilon c(\sigma))\rho_*(\sigma),
 \qquad \epsilon^2=0,
\]
has the same trace as $\rho_*$ on those elements.  Indeed, the quotient
\[
 M/(\operatorname{Ad}(\rho_*(\sigma))-1)M
\]
is the diagonal subspace in an eigenbasis, and zero evaluation kills
both diagonal entries.  Continuity and the
density just observed give equality of traces on all of $H$.

Choose $h_1,\ldots,h_4\in H$ such that $A_j=\rho_*(h_j)$ form a
basis of $M_2(K)$.  The lifts $\rho_c(h_j)$ form a basis over
$K[\epsilon]$.  The Gram matrix $\tr(A_jA_k)$ is invertible, and
equality of traces on products of two and three group elements shows
that this Gram matrix and all multiplication structure constants are
unchanged after lifting.  The linear map taking $A_j$ to
$\rho_c(h_j)$ is therefore an algebra automorphism of $M_2(K[\epsilon])$,
reducing to the identity.  Such an automorphism is conjugation by a
matrix congruent to the identity: this follows, for example, by
carrying a system of matrix units to its image and choosing the
corresponding basis.  Pairing with the four basis matrices then shows
that this conjugation carries $\rho_*|_H$ to $\rho_c|_H$ on every
element of $H$.  Hence $[c]|_H=0$.

The invariants $M^H$ are the scalars.  Inflation--restriction bounds
the kernel of all these evaluations by
\[
 H^1(\Gal(F_\infty/F),K).
\]
This space has dimension at most two.  To see the last assertion,
the corresponding tower over $\QQ$ is pro-two and unramified away
from $2$.  Its quadratic subfields are
$\QQ(i)$, $\QQ(\sqrt2)$, and $\QQ(\sqrt{-2})$.
Its elementary two-character rank is therefore two.
The field $F$ is disjoint from that tower: a nontrivial
intersection would have a quadratic subfield, whereas none of these
three fields has $2$ split completely.  Thus the tower group over
$F$ has the same two-character rank.

The ambient cohomology space is finite dimensional.  A finite list of
regular elements consequently detects its quotient by a subspace of
dimension at most two.  Extend $K$ once to split the eigenvalues of
this list.  Both regular split eigenvalues and the nonvanishing minors
of the evaluation matrix persist on sufficiently small neighborhoods.
For each $n$, Chebotarev approximation supplies Frobenius elements in
these neighborhoods which are trivial at the required finite stage
of the tower \cite[V, Theorem~3.23]{MilneCFT}.
All conditions concern finite quotients after the
neighborhoods have been chosen.  Extra such primes can be added, so
the cardinality can be any fixed sufficiently large $h$.  The dual
Selmer dimension is then at most two.

The auxiliary evaluations have reduced the undetected dual Selmer
space to at most two scalar tower directions. The next count turns
this into the uniform presentation bound \(g=2h+1\).
For completeness, write $h^j_v=\dim_K H^j(G_{F_v},M)$.
Local duality and Euler characteristics give
\[
 h^1_v=
 \begin{cases}
 2h^0_v+4[F_v:\QQ_2],&v\mid2,\\
 2h^0_v-4,&v\text{ real},\\
 2h^0_v,&v\nmid2\infty.
 \end{cases}
\]
For trivial complex conjugation, the dimensions are
$(h^0_v,h^1_v)=(4,4)$.
For nontrivial unipotent complex conjugation they are $(2,0)$.
Thus the dyadic contribution $4d$ cancels the real contribution
$-4d$.  The localization image in the direct sum of the local
$H^1$ spaces is its own annihilator under Poitou--Tate.
For a subspace of local conditions $L$, elementary linear algebra
therefore gives the Selmer difference as
$\dim L-\tfrac12\sum_v h^1_v$.
Here $L$ is zero at $S$ and full at $Q_n$.
At each auxiliary prime regularity gives $h^0_v=2$ and $h^1_v=4$.
Consequently the identity is
\[
 h^1_{\mathrm{Sel}}(M)-h^1_{\mathrm{Sel}^{\perp}}(M)
   =-\sum_{v\in S}h^0_v+2h.
\]
The finite-coefficient formulas are
\cite[I, Corollary~2.3, Theorems~2.8, 2.13, 4.10, and~5.1]{Milne2006};
their passage to $K$ is the lattice-and-limit argument of
Lemma~\ref{lem:curve-cohomology}.
With independent $S$-frames the relative tangent dimension is
\[
 h^1_{\mathrm{Sel}}(M)+\sum_{v\in S}h^0_v-h^0(G_{F,S},M)
 =h^1_{\mathrm{Sel}^{\perp}}(M)+2h-1
 \leq 2h+1.
\]
\emph{Generators with a fixed old composite.}
To obtain the last assertion as well as the variable bound, write
\(A_\eta=R_{\mathrm{loc},\eta}\) and
\(I_{\eta,n}=\ker(R_{\eta,n}\to R_\eta)\). The relative cotangent
spaces are
\[
 Q_{\eta,n}=\mathfrak m_{\eta,n}/
  (\mathfrak m_{\eta,n}^2+\mathfrak m_{A_\eta}R_{\eta,n}),
 \qquad
 Q_\eta=\mathfrak m_\eta/
  (\mathfrak m_\eta^2+\mathfrak m_{A_\eta}R_\eta).
\]
The first maps onto the second. Its kernel is the image of the
actual ideal \(I_{\eta,n}\): if a representative maps into
\(\mathfrak m_\eta^2+\mathfrak m_{A_\eta}R_\eta\), lift a
decomposition of its image into
\(\mathfrak m_{\eta,n}^2+\mathfrak m_{A_\eta}R_{\eta,n}\) and
subtract it. This leaves a representative in \(I_{\eta,n}\) of the
same relative cotangent class.

Choose a basis of the old space \(Q_\eta\) with fixed representatives
\(r_{\eta,1},\ldots,r_{\eta,q}\) in \(R_\eta\). Choose their
reductions first in the common special-fiber diagram and then lift
to both values of \(\eta\). Their relative cotangent spaces agree
modulo \(\varpi\), so the number \(q\) is the same. For each \(n\),
lift these fixed old representatives to \(R_{\eta,n}\) and complete
their independent cotangent classes to a basis using representatives
in \(I_{\eta,n}\). The tangent bound just proved shows that at most
\(g=2h+1\) elements are needed; pad the list with zeros to length
\(g\).

These choices can have identical reductions for \(\eta=1,i\).
First choose the lifts and the kernel-complement elements in the
common special-fiber diagram. Lift an element arbitrarily to
\(R_{\eta,n}\). If its image in \(R_\eta\) differs from the required
fixed old value by \(\varpi a\), subtract \(\varpi\) times any
preimage of \(a\) in \(R_{\eta,n}\). This preserves the chosen
special-fiber reduction and makes the old image exact. For a
kernel-complement element the required old image is zero. No
flatness of the old global ring over the coefficient DVR is used.

Complete relative Nakayama now gives
\eqref{eq:patching-presentation}. Every composite to \(R_\eta\)
sends the first \(q\) variables to the fixed \(r_{\eta,j}\) and
all remaining variables to zero, independently of \(n\). The two
presentations still agree modulo \(\varpi\). The variable count,
and hence every later ambient dimension, is unchanged.
\end{proof}

\begin{lemma}[Rank of auxiliary character twists]\label{lem:twist-rank}
Increase $h$ so that \mbox{$b=h-2d+2>0$}. There are
quotients
\[
 A_n=(\ZZ/2^{a_n}\ZZ)^b,\qquad a_n\longrightarrow\infty,
\]
of the ray class groups defining characters trivial at every dyadic
and real place and unramified outside $Q_n$.  At $Q_n$ they use only
the residue-unit quotients of order $2^n$.
\end{lemma}

\begin{proof}
Let $S_2$ be the dyadic places.  Take the ordinary ray class group
of modulus $\prod_{q\in Q_n}q$, and quotient by the dyadic ideal
classes.  Omitting real places from the modulus requires them to
split; at the dyadic places, both the units and uniformizers are
trivial in this quotient.  The ray class exact sequence is
\[
 \OO_{F,S_2}^{\times}\longrightarrow
 \prod_{q\in Q_n}\kk(q)^\times\longrightarrow
 \Gamma_n\longrightarrow\operatorname{Cl}_{S_2}(F)
 \longrightarrow0 .
\]
This is the usual ideal-theoretic ray sequence, with the primes
in $S_2$ inverted \cite[V, Theorem~1.7]{MilneCFT}.
The group of dyadic $S$-units has free rank
\[
 (d-1)+d=2d-1.
\]
Here we use the $S$-unit theorem
\cite[Theorem~5.11]{MilneANT}.
The direction generated by $2$ is primitive in its free part: its
valuation at each dyadic prime is one, since $2$ splits completely
in $F$.  Since $F$ is totally real, generators for these units may
therefore be chosen as $-1,2$ and $2d-2$ further units.
For $n\geq2$, put $a_n=n-1$ and take the cyclic residue-unit
quotient of order $2^{a_n}$ at each auxiliary prime.
Both $-1$ and $2$ vanish in these quotients: the first is a
$2^{a_n}$th power because $\mu_{2^{a_n+1}}$ is present, and the
second is such a power by the imposed Kummer splitting.
The elementary divisor theorem shows that quotienting
$(\ZZ/2^{a_n}\ZZ)^h$ by the images of the remaining $2d-2$
generators leaves a quotient
$(\ZZ/2^{a_n}\ZZ)^{h-2d+2}$.
The image of the local units in $\Gamma_n$ thus has such a
quotient.  A finite abelian group containing a subgroup with
this quotient also has this quotient: in invariant-factor form,
the subgroup forces at least this many factors of order at least
$2^{a_n}$ in the containing group.  Apply this to $\Gamma_n$.
Class field theory turns its quotient into the required
characters \cite[V, Theorems~5.3 and~5.5]{MilneCFT}.
Discarding the first finitely many indices is harmless.
\end{proof}

The resulting \(b=h-2d+2\) character parameters will supply the
support dimensions missing from the fixed-determinant estimate.
Their vanishing at the old local conditions is what permits twisting
without changing those conditions. Let $H_n$ denote the formal
character group of $A_n$, completed at
the trivial character and $\varpi$.  Its coordinate ring is
\[
 \mathcal L[[l_1,\ldots,l_b]]/
 \bigl((1+l_j)^{2^{a_n}}-1:1\leq j\leq b\bigr).
\]
Its tautological Galois character will be denoted $\theta_n$.
Twisting a fixed-determinant representation by $\theta_n$ preserves
the prescribed determinants at the dyadic and real places and on
bad inertia.  It therefore defines a map from the variable global
problem to the fixed global problem with these character parameters.

The local product $R_{\mathrm{loc},\eta}$ in
\eqref{eq:patching-presentation} is obtained from the enhanced
real--dyadic base $A^{\mathrm{rd}}$ of
Proposition~\ref{prop:local-geometry} by adjoining the $u$ raw tame
factors of \eqref{eq:tame-model}, with $\det\phi_j$ unrestricted.
Adjoining the $g=2h+1$ presentation variables gives
the rings $B_\eta$ of that proposition. They therefore have dimension
\begin{equation}\label{eq:patching-ambient-dimension}
 e=1+8d+4u+g=8d+4u+2h+2.
\end{equation}
Their special fibers are identified.  The underlying spectrum of
$B_i$ is irreducible and nonvertical.  All components of $B_1$
are horizontal of dimension $e$, and correspond bijectively to
the top-dimensional components of its special fiber, of dimension
$e-1$.  The generic point of each such special-fiber component lies
on just its corresponding component of $B_1$.  These precise
component assertions, and not only the dimension in
\eqref{eq:patching-ambient-dimension}, will be used below.

\subsection{The modules and their finite-level actions}

Let $N_{\eta,n}$ be the block module of completed forms at the
auxiliary ratio level $U_1(Q_n)$, with tame type
$(\eta,\eta^{-1})$ at the bad places.  Add $z$, make the enhanced
completion, and pull back to the independent $S$-frame space.
Define $N_\eta$ in the same way at old level.  The modules are
finite over $R_{\eta,n}^f$ and $R_\eta^f$, respectively, by
Propositions~\ref{prop:completed-forms} and~\ref{prop:block-comparison}.
We also regard them as modules for the variable-determinant rings
through their fixed-determinant quotients.

Choose the quotient of the diamond group
$\Delta_n=(\ZZ/2^n\ZZ)^h$, with generators $U_1,\ldots,U_h$,
and write $V_j=U_j-1$.  Lemma~\ref{lem:auxiliary-level} gives
\begin{align}
 &N_{\eta,n}\text{ is flat over }\mathcal L[\Delta_n],
                                                        \label{eq:auxiliary-flatness}\\
 &N_{\eta,n}/(V_1,\ldots,V_h)N_{\eta,n}
       \simeq N_\eta^{\oplus 2^h},                       \label{eq:old-level-action}\\
 &\overline N:=N_1/\varpi N_1\simeq N_i/\varpi N_i\ne0,
 \qquad \dim\Supp\overline N\geq 2d+4s-1.             \label{eq:old-residue-dimension}
\end{align}
The isomorphism \eqref{eq:old-level-action} is an isomorphism of
$R_{\eta,n}$-modules through $R_{\eta,n}\twoheadrightarrow R_\eta$.
In particular, it identifies the actual Galois-ring actions.
The $n$-level modules and their ring and diamond actions for the two
values of $\eta$ likewise agree modulo $\varpi$.

We recall why the operations just performed preserve the assertions.
Flatness over $\OO[\Delta_n]$ holds in the definite function-space
model and passes through the exact block functor.  The $z$-adic
flatness criterion applies after adjoining $z$, and base change
from $\OO[[z]]$ to $\mathcal L$ preserves flatness.
The subsequent localizations, noetherian completions, and smooth
frame extensions are flat maps of the commuting deformation rings.
Tensoring a finite-ideal flatness test over $\mathcal L[\Delta_n]$
with these flat extensions preserves injectivity, proving
\eqref{eq:auxiliary-flatness} after all the stated operations.
The degeneracy maps giving
\eqref{eq:old-level-action} are Hecke equivariant.  At the regular
Frobenius curve their composite is invertible, and the special
contribution is absent.  The augmented Galois action is consequently
unramified at $Q_n$ and factors through the old quotient; the
classical comparison identifies it with the old action.  Finally,
Proposition~\ref{prop:hecke-dimension} gives dimension at least
$1+2d$ on each Hecke component.  Enhanced completion preserves this
dimension, independent $S$-frames add $4s-1$, and reduction by
$\varpi$ loses at most one.  This proves the bound in
\eqref{eq:old-residue-dimension}.  Nonvanishing follows from the
pro-modular curve and the faithfulness of the old block module.

For chosen tame inertia elements $\tau_{j,n}$ at the auxiliary
primes, the commuting diamond actions also satisfy
\begin{equation}\label{eq:diamond-quadratic}
 U_j^2-T(\tau_{j,n})U_j+1=0.
\end{equation}
Here $T$ denotes the trace of the Galois deformation acting on the
module.  This is an identity of operators, as supplied by
Lemma~\ref{lem:auxiliary-level}.

We will need actual twist isomorphisms, rather than only a numerical
invariance of support.  For every sign tuple
$\sigma\in\{1,-1\}^b$, specialize $\theta_n$ to the quadratic
character $\theta_{\sigma,n}$.  Multiplication of forms by
$\theta_{\sigma,n}\circ\operatorname{Nrd}$ gives a semilinear
isomorphism of $N_{\eta,n}$ for the corresponding automorphism of
$R_{\eta,n}^f$.  The central character is unchanged because this
character is quadratic.  It is trivial above $2$ and at infinity
and unramified on bad inertia.  It preserves the ratio level: on
that level the two diagonal character values agree, so their
product is one.  Its reduction is one, hence the enhanced curve
and its residual frames are fixed.  These observations prove that
the isomorphisms act on precisely the completed modules under
consideration.

\subsection{Exact patching with finite jets}

Fix a nonprincipal ultrafilter $\mathfrak U$ on the positive integers.
For a fixed complete noetherian local ring $(B,\mathfrak m)$, define
\[
 B^*=\varprojlim_j\prod_{\mathfrak U} B/\mathfrak m^j.
\]
Its residue field is the ultrapower $K^*$ of the residue field $K$.
This deliberate auxiliary scalar extension is separate from the
finite coefficient extensions used to define the original curves
and representations.
For a fixed finite $B$-module $M$, use the same definition with
$M/\mathfrak m^jM$.  For varying quotients of $B$ and varying finite
modules, a subscript $\infty$ will denote this construction using
the maximal-ideal filtration induced by $B$.

\begin{lemma}[Finite jets and uniform exactness]\label{lem:jet-exactness}
The following assertions hold.
\begin{enumerate}
\item The ring $B^*$ is complete noetherian local, and $B\to B^*$
is faithfully flat.  For fixed finite modules the construction is
exact and equals extension of scalars.  It preserves their support
dimensions.  If $B$ is a domain, so is $B^*$.
\item For varying quotients of $B$ and modules generated by a bounded
number of elements, the patched objects are finite over $B^*$.
Surjections and cokernels of a fixed finite number of compatible
operators commute with patching.
\item Let $(D,\mathfrak a)$ be a fixed noetherian local ring and $C$
a fixed finite $D$-module.  For arbitrary finite-module extensions
\[
 0\longrightarrow P_n\longrightarrow H_n\longrightarrow C
 \longrightarrow0
\]
there is a constant $c$, independent of the extension, such that
\begin{equation}\label{eq:uniform-artin-rees}
 P_n\cap\mathfrak a^jH_n\subseteq\mathfrak a^{j-c}P_n
 \quad(j\geq c).
\end{equation}
Thus a filtration of fixed length with fixed successive quotients
patches exactly whenever its ambient jet filtration is uniformly
cofinal with the $\mathfrak a$-adic filtrations on its pieces.
\end{enumerate}
\end{lemma}

\begin{proof}
At each depth, taking an ultrapower extends the scalars in each
finite-dimensional associated graded piece from $K$ to $K^*$.
Thus the graded ring of $B^*$ is generated by the images of any
fixed finite set of generators of $\mathfrak m$.  Successive
expansion in these generators, followed by completeness, identifies
the jet filtration with powers of the ideal they generate.
The complete filtered-ring argument then proves noetherianity:
initial ideals are finitely generated in the noetherian associated
graded ring, and successive approximation lifts those generators.

For an exact sequence of fixed finite modules, Artin--Rees makes
the induced and intrinsic filtrations on its submodule cofinal.
Use the induced filtration to obtain exact sequences of jets;
cofinality identifies their inverse limits with the construction
using intrinsic filtrations.  The resulting limit is exact.  A fixed
finite presentation identifies $M^*$ with $B^*\otimes_B M$.
This gives flatness by the finite-ideal criterion, and the nonzero
residue extension gives faithful flatness.  Lengths of the graded
pieces are unchanged after residue-field extension, so their
Hilbert growth, and hence the support dimension, is unchanged.

Suppose next that $B$ is a domain and is not a field.  Choose a
discrete valuation $v$ of its fraction field centered at
$\mathfrak m$, for example one supplied by the normalized blowup
of $\mathfrak m$.  It is positive on $\mathfrak m$.  The ideals
\[
 I_a=\{x\in B:v(x)\geq a\}
\]
have zero intersection.  Chevalley's lemma for a complete local
ring implies that for every $q$ there is an $a(q)$ with
$I_{a(q)}\subseteq\mathfrak m^q$.  Thus being nonzero modulo a
fixed $\mathfrak m^q$ bounds $v(x)$ from above, uniformly in $x$.
If $c_0>0$ is the minimum valuation of fixed generators of
$\mathfrak m$, then $v(\mathfrak m^r)\geq rc_0$.
For two fixed nonzero jets, these inequalities give a depth at
which the product of any representatives of the two jets is
nonzero.  Applied on a set belonging to $\mathfrak U$, they show
that the product of two nonzero elements of $B^*$ cannot vanish.
The assertion is immediate when $B$ is a field.

For the varying construction, choose a uniform number of generators
at each level.  Each jet ultraproduct is then a quotient of a fixed
finite free module over the corresponding jet of $B^*$ and has
finite length.  Images of kernels at any fixed depth form a
descending chain in a finite-length module, hence stabilize.
The Mittag--Leffler argument gives surjectivity on inverse limits
and a finite generating set over $B^*$.  The same argument applied
to a right-exact sequence proves the cokernel assertion.  This
does not assert exactness for arbitrary varying submodules.

For the uniform assertion, choose one finite free presentation
$F\twoheadrightarrow C$, with kernel $L$, and lift it to
$f_n:F\to H_n$.  Then $H_n=P_n+f_n(F)$.  If
$x\in P_n\cap\mathfrak a^jH_n$, write
\[
 x=p+f_n(y),\qquad p\in\mathfrak a^jP_n,
 \quad y\in\mathfrak a^jF.
\]
Since $x$ maps to zero in $C$, we have $y\in L$.
Artin--Rees for the one fixed inclusion $L\subset F$ gives
$y\in\mathfrak a^{j-c}L$, with $c$ independent of $n$.
As $f_n(L)\subset P_n$, this proves
\eqref{eq:uniform-artin-rees}.  Applying this successively to
the fixed quotients of a finite filtration gives uniform
strictness for each inclusion.  Uniform cofinality transfers
these bounds to the ambient jets, proving the final statement.
\end{proof}

We now use the jet lemma in three specific ways: scalar extension
of fixed modules, bounded-generation patching of the auxiliary
modules, and uniform strictness of filtrations with fixed old
quotients. It is the fixed quotient in the third assertion, not
arbitrary varying-module exactness, that preserves the regular
sequence below.

The domain assertion in Lemma~\ref{lem:jet-exactness} can be applied
to each quotient by a minimal prime.  Exactness and nilpotence of
the nilradical then show that a fixed finite list of irreducible
components is preserved by starring.  The same applies to the
special fiber and to inclusions between its components and the
ambient components.  In particular all the geometric assertions
following \eqref{eq:patching-ambient-dimension} remain valid for
$B_\eta^*$.  This observation does not require $B_i$ itself to be
reduced.

Apply the construction to the surjections
$B_\eta\twoheadrightarrow R_{\eta,n}$ and their fixed-determinant
quotients.  The module generator numbers are uniformly bounded.
Indeed modulo the maximal ideal of $R_{\eta,n}$,
\eqref{eq:diamond-quadratic} gives $V_j^2=0$.
The quotient by all the $V_j$ is the fixed old module in
\eqref{eq:old-level-action}.  The monomials
$V_1^{e_1}\cdots V_h^{e_h}$ with $0\leq e_j\leq1$ therefore
give a uniform bound, and Nakayama supplies the required generators.
We obtain finite $N_{\eta,\infty}$ and surjections
\[
 B_\eta^*\twoheadrightarrow R_{\eta,\infty}
 \twoheadrightarrow R_{\eta,\infty}^f
 \twoheadrightarrow(R_\eta^f)^*.
\]
The commuting $U_j$ patch as well.  They satisfy the patched
versions of \eqref{eq:diamond-quadratic}; the joint action algebra
is consequently module finite over $R_{\eta,\infty}$.
Modulo its original maximal ideal all the $V_j$ are nilpotent,
so adjoining their actions does not change the complete separated
module topology.

\begin{proposition}[The patched regular sequence]
\label{prop:patched-regular-sequence}
On $N_{\eta,\infty}$ the sequence
$\varpi,V_1,\ldots,V_h$ is regular.  Moreover
\begin{align}
 \dim\Supp_{R_{\eta,\infty}}N_{\eta,\infty}
       &\geq2d+4s+h,                                  \label{eq:patched-dimension}\\
 N_{\eta,\infty}/(V_1,\ldots,V_h)N_{\eta,\infty}
       &\simeq (N_\eta^{\oplus2^h})^*.                \label{eq:patched-augmentation}
\end{align}
\end{proposition}

\begin{proof}
\emph{Polynomial layers with a fixed old action.}
Fix $\eta$, and write $\mathfrak m$ for the maximal ideal of
$B_\eta$. The presentations of Lemma~\ref{lem:auxiliary-primes}
give surjections $B_\eta\to R_{\eta,n}$ whose composites to
$R_\eta$ are independent of $n$. Regard
\[
 C=\overline N^{\oplus2^h}
\]
as a fixed finite $B_\eta$-module through this common composite.
Let $J=(\varpi,V_1,\ldots,V_h)$ denote the ideal of commuting
operators on the auxiliary modules and on their patched limit.

Fix a truncation length $r$. For all sufficiently large $n$, the
relations $(1+V_j)^{2^n}-1$ have $J$-order greater than $r$.
Flatness in \eqref{eq:auxiliary-flatness}, together with the actual
old-action identification \eqref{eq:old-level-action}, gives
multiplication isomorphisms
\begin{equation}\label{eq:finite-polynomial-layers}
 \operatorname{gr}_J^{<r} N_{\eta,n}
   \simeq C\otimes_K K[T_0,T_1,\ldots,T_h]_{<r},
\end{equation}
where $T_0,T_j$ represent $\varpi,V_j$. These are isomorphisms of
$B_\eta$-modules. Indeed the Galois action commutes with the
$J$-operators, so its action on each layer is obtained from its
action on $N_{\eta,n}/JN_{\eta,n}=C$. The fixed old composite
therefore makes every degree a fixed finite direct sum of $C$.
Only the extensions joining these layers vary with $n$.

\emph{Exactness at each fixed truncation length.}
Put $H_n=N_{\eta,n}/J^rN_{\eta,n}$ and
$F_j=J^jH_n$ for $0\leq j\leq r$. All these modules are finite
over $B_\eta$. Their successive quotients are the fixed modules
in \eqref{eq:finite-polynomial-layers}, so their numbers of
generators are bounded independently of $n$ for this fixed $r$.
If $\mathfrak q_n$ is the maximal ideal of $R_{\eta,n}$, the
surjection from $B_\eta$ gives
\[
 \mathfrak m R_{\eta,n}=\mathfrak q_n.
\]
Thus the $\mathfrak m$-adic filtrations are precisely the
maximal-ideal filtrations used for patching these modules.

Apply Lemma~\ref{lem:jet-exactness}(3) over the one fixed ring
$B_\eta$ to
\[
 0\longrightarrow F_{j+1}\longrightarrow F_j
 \longrightarrow F_j/F_{j+1}\longrightarrow0.
\]
Its fixed-quotient Artin--Rees bound makes each inclusion uniformly
strict. Composing the finitely many bounds compares the induced
filtration from $H_n$ with the intrinsic filtration on every
$F_j$. The constants may depend on $r$, but are independent of
$n$. Consequently this finite filtration patches exactly. No
exactness for arbitrary varying submodules is required.

The ideal $J^r$ is generated by a fixed finite number of monomials
in its commuting generators. The cokernel assertion of
Lemma~\ref{lem:jet-exactness}(2) therefore identifies the patched
$H_n$ with $N_{\eta,\infty}/J^rN_{\eta,\infty}$; the same
argument identifies its quotients by $J^j$. These identifications
respect the multiplication maps and are compatible as $r$ varies.
We obtain
\begin{equation}\label{eq:patched-polynomial-layers}
 \operatorname{gr}_J N_{\eta,\infty}
       \simeq C^*\otimes_{K^*}K^*[T_0,T_1,\ldots,T_h].
\end{equation}
Here $C^*$ is the scalar extension of the fixed $B_\eta$-module
$C$ from Lemma~\ref{lem:jet-exactness}(1).

\emph{Regularity, dimension, and augmentation.}
The finite joint action algebra generated by the $U_j$ is complete
local: modulo the deformation-ring maximal ideal, every $V_j$ is
nilpotent. Its module topology agrees with the one used for
patching. Since $J$ lies in its maximal ideal, the finite module
$N_{\eta,\infty}$ is $J$-adically separated and complete.
A nonzero element killed by $\varpi$ would therefore have a first
nonzero $J$-initial form, contradicting injectivity of $T_0$ in
\eqref{eq:patched-polynomial-layers}. The same initial-form
argument gives
\[
 \varpi N_{\eta,\infty}\cap J^mN_{\eta,\infty}
       =\varpi J^{m-1}N_{\eta,\infty}\qquad(m\geq1).
\]
Hence the associated graded module after quotienting by $\varpi$
is obtained by setting $T_0=0$. Repeating the argument for
$V_1,\ldots,V_h$ proves that the stated sequence is regular.

The final quotient by this regular sequence is $C^*$. Its dimension
is at least $2d+4s-1$ by \eqref{eq:old-residue-dimension} and
Lemma~\ref{lem:jet-exactness}(1). Each member of the regular
sequence raises support dimension by one on passing back to the
original module. Dimensions over the finite joint action algebra
and over $R_{\eta,\infty}$ agree, giving
\eqref{eq:patched-dimension}.
Finally, apply cokernel patching to the diamond augmentation
\eqref{eq:old-level-action}, without reducing by $\varpi$.
The old $B_\eta$-action is fixed by the same common composite,
so this gives \eqref{eq:patched-augmentation} with that actual
old deformation action.
\end{proof}

\subsection{Twisting and full support}

\begin{proposition}[Full support after patching]
\label{prop:patched-support}
The untwisted patched module satisfies
\[
 \Supp_{R_{1,\infty}}N_{1,\infty}
       \supseteq\Spec R_{1,\infty}^f.
\]
Consequently the old module has full support:
$\Supp_{R_1^f}N_1=\Spec R_1^f$.
\end{proposition}

\begin{proof}
\emph{Finite twisting and its dimension contribution.}
The finite characters of Lemma~\ref{lem:twist-rank} give a local
ring map after patching,
\begin{equation}\label{eq:formal-twisting-map}
 \Theta_\eta:R_{\eta,\infty}\longrightarrow
          R_{\eta,\infty}^f[[l_1,\ldots,l_b]].
\end{equation}
Here one takes bounded total jets in the ring and character
variables.  At each such depth the equations
$(1+l_j)^{2^{a_n}}-1$ disappear for large $n$.
Coefficientwise patching of the finitely many monomials at
each depth identifies the target with the displayed power-series
ring.  This explains both existence of the map and its target.

The map \eqref{eq:formal-twisting-map} is finite.  Since $A_n$
is a character-group quotient, choose, for each $n$ and each $j$,
a Galois element detecting exactly its $j$th coordinate.  Its
determinant divided by $\chi$ maps under twisting to
$(1+l_j)^2$.  Although these Galois elements may vary with $n$,
their determinant functions determine elements in every patched
jet, and hence in $R_{\eta,\infty}$.  The resulting elements
have residue one.  On the closed fiber of
\eqref{eq:formal-twisting-map} they impose $l_j^2=0$.
At $l_1=\cdots=l_b=0$ the map is the fixed-determinant
surjection.  Its closed fiber is therefore an Artinian local
algebra, of length at most $2^b$, with the same residue field.
Complete Nakayama proves finiteness.

Let $T_\eta$ denote
$N_{\eta,\infty}[[l_1,\ldots,l_b]]$, restricted along
\eqref{eq:formal-twisting-map} and then regarded as a finite
$B_\eta^*$-module.  Its support is the closed image of the
power-series-extended support.  Finiteness of the map preserves
its dimension, and \eqref{eq:patched-dimension} gives
\[
 \dim\Supp_{B_\eta^*}T_\eta
 \geq 2d+4s+h+b
 =8d+4u+2h+2=e.
\]
Furthermore $\varpi$ is regular on $T_\eta$, since it is
regular before power-series extension and restriction of scalars.
No minimal support prime of $T_\eta$ is vertical.

\emph{Full support on the \(i\)-model and transfer to the \(1\)-model.}
The dimension estimate now reaches the entire ambient dimension after
twisting. The spectrum of $B_i^*$ is irreducible of dimension $e$.
It follows that $T_i$ has full support.  On reducing modulo
$\varpi$, the finite-stage modules, actions, and twisting maps
for $\eta=1,i$ are identical.  Reduction is a cokernel, and
power-series extension commutes with it and with the jets.
Thus $T_1/\varpi T_1$ and $T_i/\varpi T_i$ are the same
module over the identified special fibers.  For any finite
module $T$, local Nakayama gives
\[
 \Supp(T/\varpi T)=\Supp(T)\cap V(\varpi).
\]
It follows that $\Supp T_1$ contains the generic point of
every top-dimensional special-fiber component of $B_1^*$.

We give the component argument explicitly.  Let $\mathfrak q$
be one of these generic points, and let $\mathfrak p$ be the
unique minimal prime of $B_1^*$ below it.  The local geometry
gives $\dim B_1^*/\mathfrak p=e$ and
$\dim B_1^*/\mathfrak q=e-1$; equivalently
$\mathfrak q/\mathfrak p$ has height one.  The latter also
follows from the principal ideal theorem, since it is a
minimal prime of the divisor $\varpi$ on this component.
Choose a minimal support prime $\mathfrak a\subseteq\mathfrak q$
of $T_1$.  The uniqueness assertion forces
$\mathfrak p\subseteq\mathfrak a$.  The prime $\mathfrak a$
is horizontal, so it is not $\mathfrak q$; there is no prime
strictly between $\mathfrak p$ and $\mathfrak q$.
Thus $\mathfrak a=\mathfrak p$.  Every ambient component is
obtained this way, and $T_1$ has full support.

\emph{Removing twists by actual sign intertwiners.}
It remains to remove the formal twists.  For each fixed sign
tuple $\sigma\in\{1,-1\}^b$, the finite-level quadratic
twists described above give local ring automorphisms and
semilinear module isomorphisms.  They preserve maximal-ideal
filtrations, and therefore patch, together with their inverses.
In particular the support of $N_{1,\infty}$ is stable under
the resulting sign automorphisms of $R_{1,\infty}^f$.
Evaluation of \eqref{eq:formal-twisting-map} at
$l_j=\sigma_j-1$ is the fixed-determinant quotient followed by
this patched sign twist, by its finite-level construction.
The substitution is well-defined
on formal series, since its values are $0$ or $-2$.

Let $x\in\Spec R_{1,\infty}^f$.  Full support of $T_1$
supplies a point $y$ in the support of
$N_{1,\infty}[[l_1,\ldots,l_b]]$ mapping to $x$.
The determinant calculation, since $x$ has fixed determinant,
gives in the residue field $\kappa(y)$
\[
 (1+l_j)^2=1,\qquad\text{or equivalently}\qquad l_j(l_j+2)=0.
\]
For every $j$, primality therefore chooses $l_j=0$ or $-2$.
Choose the corresponding sign tuple.  Quotienting by these
sign equations leaves a nonzero localized module, by Nakayama,
and identifies the quotient with $N_{1,\infty}$ at a point
$x'$ of its support.  Its image $x$ is the sign transform of
$x'$.  Sign stability puts $x$ in the untwisted support.
In characteristic two the two signs coincide, with the same
conclusion.  This proves the first assertion.

\emph{Diamond augmentation and the old action.}
Finally consider the old quotient $(R_1^f)^*$.  Let
$\mathfrak p$ be any of its primes, viewed in
$\Spec R_{1,\infty}$.  It belongs to the support just proved.
In the finite joint action algebra generated by the $U_j$,
there is a support prime $\mathfrak P$ above $\mathfrak p$.
On the old quotient each $\tau_{j,n}$ is trivial, so
\eqref{eq:diamond-quadratic} becomes
\[
 U_j^2-2U_j+1=V_j^2=0.
\]
Hence $V_j\in\mathfrak P$ for all $j$.  The nonzero finite
module localized at $\mathfrak P$ remains nonzero after
quotienting by all the $V_j$, by Nakayama.  Equation
\eqref{eq:patched-augmentation} shows that
$(N_1^{\oplus2^h})^*$ has support at every prime of
$(R_1^f)^*$.  Faithful flatness in
Lemma~\ref{lem:jet-exactness} now gives
$\Supp_{R_1^f}N_1=\Spec R_1^f$.
\end{proof}

\begin{proof}[Proof of Theorem~\ref{thm:localized-propagation}]
We have proved full support on the enhanced completed global
deformation germ with independent $S$-frames, global determinant
$\chi$, and unipotent bad inertia.  After the initial base
changes, the source locus has exactly these local conditions.
At its absolutely irreducible curve, the representation and
pseudo-deformation problems agree after a smooth choice of
frames.  Adding $z$, passing to a finite curve-field extension,
and taking this completed localization are faithfully flat on
the indicated germ, by Lemma~\ref{lem:curve-completion}.
Consequently the closed Hecke support contains that source germ.

To pass from the germ to the asserted irreducible closed locus,
work on its reduced irreducible coordinate ring.  The image of
any equation of the Hecke support is nilpotent in the completed
source germ, by the set-theoretic containment just proved.
Faithful flatness descends this nilpotence to the localization
at the enhanced curve.  That localization is reduced, so the
image is zero there.  It therefore vanishes on the entire
irreducible locus: localization of a domain is injective.
Thus the locus is pro-modular over the field reached
by the allowed solvable base changes.  This is precisely
potential pro-modularity and proves the theorem.
\end{proof}
